% !TEX program = xelatex
\documentclass[12pt,a4paper]{article}

% ---------- Font và ngôn ngữ ----------
\usepackage{fontspec}
\setmainfont{texgyretermes}[
  Extension      = .otf,
  UprightFont    = *-regular,
  BoldFont       = *-bold,
  ItalicFont     = *-italic,
  BoldItalicFont = *-bolditalic]
\setmonofont{DejaVu Sans Mono}[Scale=0.85]
\usepackage{polyglossia}
\setmainlanguage{vietnamese}

% ---------- Bố cục ----------
\usepackage[a4paper,
  left=3cm, right=3cm, top=2.5cm, bottom=2.5cm,
  headheight=15pt, headsep=0.8cm, footskip=1cm]{geometry}
\usepackage{graphicx}
\usepackage{etoolbox}
\usepackage{needspace}
\usepackage{amsmath}
\usepackage{setspace}
\setstretch{1.15}
\usepackage{parskip}
\usepackage{enumitem}
\setlist{nosep, leftmargin=*}

% ---------- Bảng, màu, hình ----------
\usepackage{array, booktabs, tabularx, longtable, multirow, makecell}
\usepackage[table]{xcolor}
\usepackage{tikz}
\usetikzlibrary{positioning, arrows.meta, fit, backgrounds, shapes.geometric, shapes.misc}
\usepackage{caption}
\captionsetup{font=small, labelfont=bf}
\usepackage{tcolorbox}
\tcbuselibrary{skins}

% ---------- Tiêu đề mục, header ----------
\definecolor{primary}{HTML}{7A1E3A}
\definecolor{accent}{HTML}{1F5FA8}
\definecolor{softgray}{HTML}{F3F3F3}
\usepackage{titlesec}
\titleformat{\section}{\Large\bfseries\color{primary}}{\thesection.}{0.5em}{}
\titleformat{\subsection}{\large\bfseries\color{accent}}{\thesubsection}{0.5em}{}
\titlespacing*{\section}{0pt}{1.2em}{0.5em}
\titlespacing*{\subsection}{0pt}{0.9em}{0.3em}

%% ------------------------------------------------------------
%% THÔNG TIN BÁO CÁO (dùng cho bìa, header, footer)
%% ------------------------------------------------------------
\newcommand{\tenDHQG}{ĐẠI HỌC QUỐC GIA THÀNH PHỐ HỒ CHÍ MINH}
\newcommand{\tenTruong}{TRƯỜNG ĐẠI HỌC KHOA HỌC TỰ NHIÊN}
\newcommand{\tenKhoa}{KHOA CÔNG NGHỆ THÔNG TIN}
\newcommand{\loaiBaoCao}{KẾ HOẠCH THỰC HIỆN NỘI BỘ}
\newcommand{\tenNganh}{Môn: Vật lý cho Công nghệ thông tin}
\newcommand{\tenDeTaiChinh}{Phòng học thông minh}
\newcommand{\tenDeTaiPhu}{Giám sát chất lượng không khí và cảnh báo nhiệt độ, độ ẩm}
\newcommand{\tenDeTaiAnh}{Smart Classroom: Air Quality Monitoring and Temperature--Humidity Alert}
\newcommand{\tenDeTaiNgan}{Kế hoạch nội bộ --- Đồ án Vật lý cho CNTT}   % hiển thị trên header
\newcommand{\giangVien}{Bùi Thanh Lâm, Cao Xuân Nam, Đặng Hoài Thương}
\newcommand{\tenNhom}{}   % để trống thì bìa không hiện dòng tên nhóm

\newcommand{\svA}{Nguyễn Gia Quốc Uy}  \newcommand{\mssvA}{24127261}
\newcommand{\svB}{Ngô Quốc Tuấn}       \newcommand{\mssvB}{24127581}
\newcommand{\svC}{Nguyễn Quốc Thành}   \newcommand{\mssvC}{24127542}

\newcommand{\khoaHoc}{2026 -- 2027}
\newcommand{\thangNam}{Tháng 10 năm 2026}
\newcommand{\diaDiem}{TP. Hồ Chí Minh}
\newcommand{\fileLogo}{logo-hcmus.png}   % đặt logo cùng thư mục (tùy chọn)

%% ------------------------------------------------------------
%% HEADER & FOOTER (theo format.tex)
%% ------------------------------------------------------------
\usepackage{fancyhdr}
\pagestyle{fancy}
\fancyhf{}
\renewcommand{\headrulewidth}{0.4pt}
\renewcommand{\footrulewidth}{0.4pt}

\fancyhead[L]{\small\textit{Trường ĐH Khoa học Tự nhiên -- ĐHQG-HCM}}
\fancyhead[R]{\small\textit{\tenDeTaiNgan}}
\fancyfoot[L]{\small\textit{\MakeUppercase{\tenKhoa}}}
\fancyfoot[R]{\small\thepage}

\fancypagestyle{plain}{%
  \fancyhf{}
  \renewcommand{\headrulewidth}{0pt}
  \renewcommand{\footrulewidth}{0.4pt}
  \fancyfoot[L]{\small\textit{\MakeUppercase{\tenKhoa}}}
  \fancyfoot[R]{\small\thepage}
}

%% ------------------------------------------------------------
%% TRANG BÌA (theo format.tex: khung viền kép)
%% ------------------------------------------------------------
\newcommand{\makecover}{%
  \begin{titlepage}
  \thispagestyle{empty}
  \singlespacing
  \begin{tikzpicture}[remember picture, overlay, inner sep=0, outer sep=0]
    % Khung chinh
    \draw[blue!55!black, line width=5.5pt]
      ([xshift=-1.5cm, yshift=-2cm]current page.north east) coordinate (A) --
      ([xshift=1.5cm, yshift=-2cm]current page.north west) coordinate (B) --
      ([xshift=1.5cm, yshift=2cm]current page.south west) coordinate (C) --
      ([xshift=-1.5cm, yshift=2cm]current page.south east) coordinate (D) -- cycle;

    % Vien long 1
    \draw
      ([yshift=0.2cm,xshift=-0.3cm]A)-- ([yshift=0.2cm,xshift=0.3cm]B)--
      ([yshift=-0.3cm,xshift=0.3cm]B) --([yshift=-0.3cm,xshift=-0.2cm]B)--
      ([yshift=0.3cm,xshift=-0.2cm]C)--([yshift=0.3cm,xshift=0.3cm]C)--
      ([yshift=-0.2cm,xshift=0.3cm]C)-- ([yshift=-0.2cm,xshift=-0.3cm]D)--
      ([yshift=0.3cm,xshift=-0.3cm]D)--([yshift=0.3cm,xshift=0.2cm]D)--
      ([yshift=-0.3cm,xshift=0.2cm]A)--([yshift=-0.3cm,xshift=-0.3cm]A)--
      ([yshift=0.2cm,xshift=-0.3cm]A);

    % Vien long 2
    \draw
      ([yshift=-0.2cm,xshift=0.2cm]A)-- ([yshift=-0.2cm,xshift=-0.2cm]B)--
      ([yshift=0.2cm,xshift=-0.2cm]B) --([yshift=0.2cm,xshift=0.2cm]B)--
      ([yshift=-0.2cm,xshift=0.2cm]C)--([yshift=-0.2cm,xshift=-0.2cm]C)--
      ([yshift=0.2cm,xshift=-0.2cm]C)-- ([yshift=0.2cm,xshift=0.2cm]D)--
      ([yshift=-0.2cm,xshift=0.2cm]D)--([yshift=-0.2cm,xshift=-0.2cm]D)--
      ([yshift=0.2cm,xshift=-0.2cm]A)--([yshift=0.2cm,xshift=0.2cm]A)--
      ([yshift=-0.2cm,xshift=0.2cm]A);
  \end{tikzpicture}
  \centering
  \vspace*{0.3cm}

  {\fontsize{14pt}{17pt}\selectfont\bfseries \tenDHQG\par}
  \vspace{2pt}
  {\fontsize{14pt}{17pt}\selectfont\bfseries \tenTruong\par}
  \vspace{2pt}
  {\fontsize{14pt}{17pt}\selectfont\bfseries \tenKhoa\par}
  \vspace{0.6cm}
  \rule{4cm}{0.8pt}

  \vspace{0.8cm}
  \IfFileExists{\fileLogo}{\includegraphics[height=3.5cm]{\fileLogo}}%
    {\fbox{\parbox[c][3.5cm][c]{3.5cm}{\centering\small [Logo HCMUS]\\(đặt file \texttt{\fileLogo})}}}

  \vspace{0.9cm}
  {\fontsize{14pt}{18pt}\selectfont\bfseries\color{black} \loaiBaoCao\par}
  \ifdefempty{\tenNganh}{}{\vspace{6pt}{\fontsize{15pt}{20pt}\selectfont\color{black} \tenNganh\par}}

  \vspace{0.75cm}
  \begin{minipage}{0.94\textwidth}
    \centering
    {\fontsize{19pt}{25pt}\selectfont\bfseries\color{black}
      \tenDeTaiChinh\par}
    \vspace{3pt}
    {\fontsize{17pt}{23pt}\selectfont\color{black}
      \tenDeTaiPhu\par}
    \vspace{0.3cm}
    {\fontsize{15pt}{20pt}\selectfont\itshape\color{black}
      (\tenDeTaiAnh)\par}
  \end{minipage}\par
  \ifdefempty{\tenNhom}{}{%
    \vspace{0.3cm}
    {\color{black}\rule{2.4cm}{0.5pt}\par}
    \vspace{0.25cm}
    {\fontsize{14pt}{18pt}\selectfont\bfseries\color{black} \tenNhom\par}}

  \vfill

  \begin{center}
  \begin{tabular}{@{}r@{\hspace{0.4cm}}l@{}}
    \textbf{Giảng viên hướng dẫn:} & Bùi Thanh Lâm \\[2pt]
                                   & Cao Xuân Nam \\[2pt]
                                   & Đặng Hoài Thương \\[4pt]
    \textbf{Sinh viên thực hiện:}  & \svA \ -- \mssvA \\[2pt]
                                   & \svB \ -- \mssvB \\[2pt]
                                   & \svC \ -- \mssvC \\[4pt]
    \textbf{Năm học:}              & \khoaHoc \\
  \end{tabular}
  \end{center}

  \vfill

  {\fontsize{13pt}{16pt}\selectfont\bfseries \diaDiem, \thangNam\par}
  \vspace{0.3cm}
  \end{titlepage}
  \setcounter{page}{1}
}

\usepackage[hidelinks]{hyperref}

\newcolumntype{L}[1]{>{\raggedright\arraybackslash}p{#1}}
\newcolumntype{C}[1]{>{\centering\arraybackslash}p{#1}}
\newcolumntype{Y}{>{\raggedright\arraybackslash}X}
% Bảng dữ liệu: viền ngoài đậm (I), đường kẻ bên trong nhạt (i)
\definecolor{gridgray}{gray}{0.7}
\arrayrulecolor{gridgray}
\newcolumntype{I}{!{{\color{black}\vrule width 0.9pt}}}
\newcolumntype{i}{!{{\color{gridgray}\vrule width 0.4pt}}}
\newcommand{\outerline}{\Xhline{0.9pt}}
\renewcommand{\arraystretch}{1.25}

\newtcolorbox{notebox}[1][]{colback=softgray, colframe=primary, boxrule=0.6pt,
  arc=2pt, left=6pt, right=6pt, top=4pt, bottom=4pt, #1}

\begin{document}

% ======================================================================
% TRANG BÌA
% ======================================================================
\makecover

\pagenumbering{roman}

\tableofcontents
\clearpage
\pagenumbering{arabic}

\noindent\begin{minipage}{\linewidth}
\setlength{\parskip}{6pt}
% ======================================================================
\section{Giới thiệu tài liệu}
% ======================================================================
Đây là tài liệu nội bộ của nhóm, đi kèm bản đề xuất (proposal). Nội dung đề tài, danh sách use case, kiến trúc, các luồng xử lý, thiết bị và chi phí đã có trong proposal nên không nhắc lại ở đây. Tài liệu này chỉ ghi: ai làm gì, các quy ước chung để code của ba người khớp với nhau, thứ tự thực hiện và các rủi ro khi làm.

Mã use case (UC1--UC11) và số luồng (luồng 1--8) dùng theo proposal.
\end{minipage}
\par\bigskip

\noindent\begin{minipage}{\linewidth}
\setlength{\parskip}{6pt}
% ======================================================================
\section{Tổng quan phân công}
% ======================================================================
\begin{tabularx}{\linewidth}{I L{1.6cm} i L{2.2cm} i L{2.6cm} i Y i L{2.3cm} I}
\outerline
\textbf{Thành viên} & \textbf{Thiết bị} & \textbf{Use case} & \textbf{Yêu cầu của giảng viên} & \textbf{Việc chung chính} \\
\hline
Thành & DHT22 & UC1, UC6, UC8 & Căn bản luồng (1); nâng cao \#5, \#9 & Tích hợp Git \\
\hline
Uy & MQ-135 & UC2, UC5, UC9; phục vụ UC11 & Căn bản luồng (1); nâng cao \#4, \#12 & Lắp ráp, đóng hộp \\
\hline
Tuấn & LED RGB, Buzzer & UC3, UC4, UC7; phục vụ UC10 & Căn bản luồng (2); nâng cao \#3, \#6 & Khung website (\#10) \\
\outerline
\end{tabularx}

Mỗi người sở hữu thiết bị của mình và đi trọn luồng của thiết bị đó qua bốn lớp: phần cứng, firmware, backend, frontend. Khi vấn đáp, mỗi người phải giải thích được toàn bộ luồng mình phụ trách.
\end{minipage}
\par\bigskip

% ======================================================================
\section{Nhiệm vụ chi tiết của từng thành viên}
% ======================================================================
Trong mỗi bảng, nhiệm vụ xếp theo thứ tự: luồng căn bản, các yêu cầu nâng cao, rồi phần phục vụ người khác. Các bước trong luồng căn bản đi theo chiều dữ liệu.

\noindent\begin{minipage}{\linewidth}
\setlength{\parskip}{6pt}
\subsection{Thành --- Nhiệt độ, độ ẩm và tài khoản}
\begin{tabularx}{\linewidth}{I C{1.1cm} i L{1.7cm} i Y i L{3.1cm} i C{1.3cm} I}
\outerline
\textbf{UC} & \textbf{Lớp} & \textbf{Nhiệm vụ} & \textbf{Yêu cầu} & \textbf{Điểm} \\
\hline
\multirow{4}{*}{UC1} & Phần cứng & Đấu nối DHT22 (điện trở kéo lên 10 k$\Omega$) & \multirow{4}{=}{Căn bản, luồng (1): DHT22 $\rightarrow$ ESP $\rightarrow$ MQTT $\rightarrow$ BE $\rightarrow$ FE} & \multirow{4}{*}{\makecell{4,0\\(chung)}} \\
\cline{2-3}
 & Firmware & Đọc DHT22 định kỳ, publish topic \texttt{dht22} & & \\
\cline{2-3}
 & Backend & Subscribe \texttt{dht22}, đẩy dữ liệu xuống frontend qua WebSocket & & \\
\cline{2-3}
 & Frontend & Hiển thị nhiệt độ, độ ẩm dạng text và gauge & & \\
\hline
UC6 & Backend, Frontend & API truy vấn lịch sử từ Firebase theo khoảng thời gian; vẽ biểu đồ nhiệt độ, độ ẩm, CO\textsubscript{2} theo giờ/ngày & Nâng cao \#5 (chart) & 1,5 \\
\hline
UC8 & Backend, Frontend & Đăng ký, đăng nhập; băm mật khẩu; cấp phiên và bảo vệ các API; gắn thiết bị với tài khoản bằng mã & Nâng cao \#9 (database) & 1,5 \\
\outerline
\end{tabularx}

\textbf{Ôn vấn đáp:} luồng 1, luồng 3 (phần biểu đồ), luồng 7; nguyên lý vật lý của DHT22.
\end{minipage}
\par\bigskip

\noindent\begin{minipage}{\linewidth}
\setlength{\parskip}{6pt}
\subsection{Uy --- CO\textsubscript{2}, lưu trữ và cấu hình}
\begin{tabularx}{\linewidth}{I C{1.1cm} i L{1.7cm} i Y i L{3.1cm} i C{1.3cm} I}
\outerline
\textbf{UC} & \textbf{Lớp} & \textbf{Nhiệm vụ} & \textbf{Yêu cầu} & \textbf{Điểm} \\
\hline
\multirow{4}{*}{UC2} & Phần cứng & Đấu nối MQ-135 qua cầu chia áp (tối đa $\approx$ 2,5 V) vào chân ADC1; đo $R_L$ thực tế trên module & \multirow{4}{=}{Căn bản, luồng (1): MQ-135 $\rightarrow$ ESP $\rightarrow$ MQTT $\rightarrow$ BE $\rightarrow$ FE} & \multirow{4}{*}{\makecell{4,0\\(chung)}} \\
\cline{2-3}
 & Firmware & Đọc ADC, tính $R_S$, hiệu chuẩn $R_0$, quy đổi ra ppm, lọc nhiễu; publish topic \texttt{mq135} & & \\
\cline{2-3}
 & Backend & Subscribe \texttt{mq135}, đẩy dữ liệu xuống frontend qua WebSocket & & \\
\cline{2-3}
 & Frontend & Hiển thị CO\textsubscript{2} dạng text và gauge & & \\
\hline
UC5 & Backend & Thiết kế cấu trúc dữ liệu Firebase; lưu dữ liệu cả hai cảm biến theo thời gian & Nâng cao \#4 (theo thời gian) & 1,5 \\
\hline
UC9 & Firmware & Khung firmware kết nối WiFi và MQTT; trang cấu hình WiFi lần đầu (WiFiManager); giữ nút BOOT 5 giây để xóa & Nâng cao \#12 (dịch vụ sẵn có) & 1,0 \\
\hline
UC11 & Backend, Frontend & Trang Cài đặt: thanh trượt chỉnh ngưỡng; lưu ngưỡng vào Firebase; cung cấp ngưỡng cho Tuấn & Phục vụ Tuấn & --- \\
\outerline
\end{tabularx}

\textbf{Ôn vấn đáp:} luồng 2, luồng 3 (phần lưu trữ), luồng 8; nguyên lý MQ-135 và phương pháp ước lượng CO\textsubscript{2} (mục 5.2 proposal).
\end{minipage}
\par\bigskip

\noindent\begin{minipage}{\linewidth}
\setlength{\parskip}{6pt}
\subsection{Tuấn --- LED RGB, còi và thông báo}
\begin{tabularx}{\linewidth}{I C{1.1cm} i L{1.7cm} i Y i L{3.1cm} i C{1.3cm} I}
\outerline
\textbf{UC} & \textbf{Lớp} & \textbf{Nhiệm vụ} & \textbf{Yêu cầu} & \textbf{Điểm} \\
\hline
\multirow{4}{*}{UC3} & Frontend & Nút chọn chế độ (thủ công / tự động) và chọn màu LED & \multirow{4}{=}{Căn bản, luồng (2): FE $\rightarrow$ BE $\rightarrow$ MQTT $\rightarrow$ ESP $\rightarrow$ LED RGB} & \multirow{4}{*}{\makecell{4,0\\(chung)}} \\
\cline{2-3}
 & Backend & Nhận lệnh; chế độ tự động lấy mức chất lượng (UC10); publish \texttt{led/cmd} & & \\
\cline{2-3}
 & Firmware & Subscribe \texttt{led/cmd}, điều khiển WS2812 & & \\
\cline{2-3}
 & Phần cứng & Đấu nối LED WS2812 (điện trở 330 $\Omega$ trên đường tín hiệu) & & \\
\hline
\multirow{4}{*}{UC4} & Frontend & Nút thử còi, tắt còi & \multirow{4}{=}{Nâng cao \#3: FE $\rightarrow$ BE $\rightarrow$ MQTT $\rightarrow$ ESP $\rightarrow$ Buzzer} & \multirow{4}{*}{1,5} \\
\cline{2-3}
 & Backend & Nhận lệnh từ nút; tự bật còi khi nhiệt độ, độ ẩm vượt ngưỡng (UC11); publish \texttt{buzzer/cmd} & & \\
\cline{2-3}
 & Firmware & Subscribe \texttt{buzzer/cmd}, bật/tắt còi & & \\
\cline{2-3}
 & Phần cứng & Đấu nối buzzer qua transistor & & \\
\hline
UC7 & Backend, Frontend & Nút ``Kết nối Telegram''; lưu mã cuộc trò chuyện; gửi thông báo khi CO\textsubscript{2} ở mức Kém liên tục 10 phút & Nâng cao \#6 & 1,0 \\
\hline
UC10 & Backend, Frontend & Quy đổi CO\textsubscript{2} ra mức Tốt / Trung bình / Kém; dùng cho LED tự động và Telegram; hiển thị mức trên website & Phục vụ UC3, UC7 & --- \\
\outerline
\end{tabularx}

\textbf{Ôn vấn đáp:} luồng 4, luồng 5, luồng 6; nguyên lý LED RGB và buzzer.
\end{minipage}
\par\bigskip

\noindent\begin{minipage}{\linewidth}
\setlength{\parskip}{6pt}
\subsection{Việc chung}
Mỗi việc chung có một người chịu trách nhiệm chính; các thành viên khác hỗ trợ khi cần.

\begin{tabularx}{\linewidth}{I Y i L{2.6cm} i C{1.6cm} i C{1.4cm} I}
\outerline
\textbf{Nhiệm vụ} & \textbf{Yêu cầu} & \textbf{Phụ trách} & \textbf{Giờ} \\
\hline
Tạo repository, quy ước nhánh Git, rà soát và hợp nhất (merge) code & Làm việc nhóm & Thành & 6 \\
\hline
Mua linh kiện; hàn mạch lên board đồng lỗ; lắp hộp, khoét lỗ thông khí & Sản phẩm thật & Uy & 6 \\
\hline
Dựng khung website (bố cục, điều hướng, giao diện chung); mỗi người gắn phần của mình vào & Nâng cao \#10 & Tuấn & 6 \\
\hline
Báo cáo: mỗi người viết phần UC của mình; phần chung chia đều & Báo cáo & Cả nhóm & 3 / người \\
\outerline
\end{tabularx}
\end{minipage}
\par\bigskip

% ======================================================================
\section{Quy ước chung giữa các thành viên}
% ======================================================================
Các quy ước dưới đây cần chốt ngay giai đoạn 1 để code của ba người ghép được với nhau.

\noindent\begin{minipage}{\linewidth}
\setlength{\parskip}{6pt}
\subsection{Topic MQTT và định dạng bản tin}
Tiền tố \texttt{<id>} là mã riêng của từng thiết bị.

\begin{tabularx}{\linewidth}{I L{2.5cm} i L{2.1cm} i Y i C{1.5cm} I}
\outerline
\textbf{Topic} & \textbf{Hướng} & \textbf{Payload ví dụ} & \textbf{Phụ trách} \\
\hline
\texttt{<id>/room1/} \texttt{dht22} & ESP32 $\rightarrow$ Backend & \texttt{\{"temp":28.5,"hum":62\}} & Thành \\
\hline
\texttt{<id>/room1/} \texttt{mq135} & ESP32 $\rightarrow$ Backend & \texttt{\{"co2\_ppm":850,"raw":1203\}} & Uy \\
\hline
\texttt{<id>/room1/} \texttt{led/cmd} & Backend $\rightarrow$ ESP32 & \texttt{\{"mode":"auto"\}} hoặc \texttt{\{"mode":"manual",} \texttt{"color":"\#00FF00"\}} & Tuấn \\
\hline
\texttt{<id>/room1/} \texttt{buzzer/cmd} & Backend $\rightarrow$ ESP32 & \texttt{\{"state":"on",} \texttt{"reason":"temp\_high"\}} hoặc \texttt{\{"state":"off"\}} & Tuấn \\
\outerline
\end{tabularx}
\end{minipage}
\par\bigskip

\noindent\begin{minipage}{\linewidth}
\setlength{\parskip}{6pt}
\subsection{Sơ đồ chân ESP32}
\begin{center}
\begin{tabular}{I L{3.0cm} i L{2.5cm} i L{5.2cm} i C{1.4cm} I}
\outerline
\textbf{Thiết bị} & \textbf{Chân ESP32} & \textbf{Ghi chú} & \textbf{Đấu nối} \\
\hline
DHT22 (data) & GPIO 4 & Điện trở kéo lên 10 k$\Omega$ & Thành \\
\hline
MQ-135 (AO) & GPIO 34 (ADC1) & Cầu chia áp 5 V $\rightarrow$ tối đa $\approx$ 2,5 V & Uy \\
\hline
LED WS2812 (DIN) & GPIO 5 & Nối tiếp điện trở 330 $\Omega$ & Tuấn \\
\hline
Buzzer & GPIO 26 & Qua transistor & Tuấn \\
\hline
Nút BOOT (có sẵn) & GPIO 0 & Giữ 5 giây để xóa cấu hình WiFi & Uy \\
\outerline
\end{tabular}
\end{center}
MQ-135 phải dùng chân ADC1, vì ADC2 của ESP32 không đọc được khi WiFi đang hoạt động.
\end{minipage}
\par\bigskip

\noindent\begin{minipage}{\linewidth}
\setlength{\parskip}{6pt}
\subsection{Cấu trúc repository và người sở hữu}
\begin{tabularx}{\linewidth}{I L{4.2cm} i Y i C{1.6cm} I}
\outerline
\textbf{Thư mục, module} & \textbf{Nội dung} & \textbf{Phụ trách} \\
\hline
\texttt{firmware/dht} & Đọc DHT22, publish \texttt{dht22} & Thành \\
\hline
\texttt{firmware/mq135}, \texttt{firmware/wifi} & Đọc MQ-135, hiệu chuẩn; kết nối WiFi, MQTT, WiFiManager & Uy \\
\hline
\texttt{firmware/led}, \texttt{firmware/buzzer} & Nhận lệnh, điều khiển LED và còi & Tuấn \\
\hline
\texttt{backend/auth}, \texttt{backend/history} & Đăng nhập, phiên, gắn thiết bị; API lịch sử cho biểu đồ & Thành \\
\hline
\texttt{backend/storage}, \texttt{backend/settings} & Ghi Firebase; lưu và đọc ngưỡng & Uy \\
\hline
\texttt{backend/quality}, \texttt{backend/control}, \texttt{backend/telegram} & Mức chất lượng; lệnh LED, còi; thông báo Telegram & Tuấn \\
\hline
\texttt{frontend/} (khung chung) & Bố cục, điều hướng; mỗi người thêm trang của mình & Tuấn \\
\outerline
\end{tabularx}
\end{minipage}
\par\bigskip

\noindent\begin{minipage}{\linewidth}
\setlength{\parskip}{6pt}
\subsection{Phụ thuộc giữa các thành viên}
\begin{tabularx}{\linewidth}{I C{1.5cm} i Y i C{1.5cm} i C{1.8cm} I}
\outerline
\textbf{Người cần} & \textbf{Cần gì} & \textbf{Người cung cấp} & \textbf{Hạn} \\
\hline
Cả nhóm & Khung firmware kết nối WiFi và MQTT để gắn module của mình vào & Uy & Giai đoạn 1 \\
\hline
Cả nhóm & Khung website để gắn giao diện của mình vào & Tuấn & Giai đoạn 1 \\
\hline
Thành & Cấu trúc dữ liệu Firebase để viết API lịch sử, vẽ biểu đồ (UC6) & Uy & Giai đoạn 1 \\
\hline
Tuấn & Ngưỡng từ trang Cài đặt cho còi tự động (UC4) và Telegram (UC7); trước đó dùng ngưỡng mặc định & Uy & Giai đoạn 4 \\
\hline
Cả nhóm & Phiên đăng nhập để bảo vệ API; trước đó phát triển không cần đăng nhập & Thành & Giai đoạn 4 \\
\outerline
\end{tabularx}
\end{minipage}
\par\bigskip

% ======================================================================
\section{Thống kê khối lượng công việc}
% ======================================================================
\noindent\begin{minipage}{\linewidth}
\setlength{\parskip}{6pt}
Số giờ dưới đây là ước lượng tương đối để so sánh giữa các thành viên, không phải cam kết thời gian.

\begin{tabularx}{\linewidth}{I C{1.7cm} i Y i C{1.6cm} i C{1.5cm} i C{1.2cm} I}
\outerline
\textbf{Thành viên} & \textbf{Use case (giờ ước lượng)} & \textbf{Chuyên môn} & \textbf{Việc chung} & \textbf{Tổng} \\
\hline
Thành & UC1 nhiệt độ, độ ẩm (8); UC6 biểu đồ (8); UC8 đăng ký, đăng nhập (12) & 28 & 9 & \textbf{37} \\
\hline
Uy & UC2 CO\textsubscript{2}, có hiệu chuẩn MQ-135 (11); UC5 lưu Firebase (6); UC9 cấu hình WiFi (5); UC11 trang Cài đặt (6) & 28 & 9 & \textbf{37} \\
\hline
Tuấn & UC3 LED RGB (10); UC4 còi (6); UC7 Telegram (8); UC10 mức chất lượng (2) & 26 & 9 & \textbf{35} \\
\outerline
\end{tabularx}
\end{minipage}
\par\bigskip

\noindent\begin{minipage}{\linewidth}
\setlength{\parskip}{6pt}
\begin{tabularx}{\linewidth}{I L{3.6cm} i Y i Y i Y I}
\outerline
\textbf{Tiêu chí} & \textbf{Thành} & \textbf{Uy} & \textbf{Tuấn} \\
\hline
Thiết bị phụ trách & DHT22 & MQ-135 & LED RGB, Buzzer \\
\hline
Lớp được làm & Phần cứng, firmware, backend, frontend & Phần cứng, firmware, backend, frontend & Phần cứng, firmware, backend, frontend \\
\hline
Số use case & 3 & 4 & 4 \\
\hline
Yêu cầu nâng cao & \#5, \#9 & \#4, \#12 & \#3, \#6 \\
\hline
Phần khó nhất & Đăng nhập, gắn thiết bị với tài khoản & Hiệu chuẩn MQ-135, cấu hình WiFi trên thiết bị thật & Hai luồng điều khiển, logic cảnh báo \\
\outerline
\end{tabularx}

\textbf{Nhận xét:} chênh lệch tổng khối lượng khoảng 2 giờ (dưới 6\%), nằm trong sai số ước lượng. Cả ba người đều làm đủ bốn lớp từ phần cứng đến giao diện.
\end{minipage}
\par\bigskip

\noindent\begin{minipage}{\linewidth}
\setlength{\parskip}{6pt}
% ======================================================================
\section{Kế hoạch thực hiện}
% ======================================================================
\begin{tabularx}{\linewidth}{I C{1.3cm} i Y i Y i Y I}
\outerline
\textbf{Giai đoạn} & \textbf{Thành} & \textbf{Uy} & \textbf{Tuấn} \\
\hline
1 & Tạo repository, quy ước nhánh Git & Mua linh kiện, làm nóng MQ-135; cấu trúc dữ liệu Firebase; khung firmware WiFi, MQTT & Khung website; chốt topic MQTT và payload \\
\hline
2 & Firmware đọc DHT22, publish & Firmware đọc MQ-135, hiệu chuẩn $R_0$, publish & Firmware nhận lệnh LED, còi \\
\hline
3 & Luồng 1 lên website (text, gauge) & Luồng 2 lên website; lưu Firebase (UC5) & Luồng 4: điều khiển LED từ web \\
\hline
4 & Biểu đồ (UC6); đăng nhập (UC8) & Cấu hình WiFi (UC9); trang Cài đặt (UC11) & Còi (UC4); mức chất lượng, LED tự động (UC10); Telegram (UC7) \\
\hline
5 & Tích hợp, hợp nhất code & Hàn mạch, đóng hộp & Kiểm thử cài đặt với người không chuyên \\
\hline
6 & Viết báo cáo phần của mình; ôn vấn đáp luồng 1, 3, 7 & Viết báo cáo phần của mình; ôn vấn đáp luồng 2, 3, 8 & Viết báo cáo phần của mình; ôn vấn đáp luồng 4, 5, 6 \\
\outerline
\end{tabularx}
\end{minipage}
\par\bigskip

\noindent\begin{minipage}{\linewidth}
\setlength{\parskip}{6pt}
% ======================================================================
\section{Dự kiến điểm theo quy định}
% ======================================================================
\begin{center}
\begin{tabular}{I L{3.9cm} i C{1.2cm} i C{2.2cm} i C{1.8cm} i L{3.0cm} I}
\outerline
\textbf{Hạng mục} & \textbf{Tối đa} & \textbf{Kế hoạch} & \textbf{Phụ trách} & \textbf{Ghi chú} \\
\hline
Chức năng căn bản & 4,0 & 4,0 & Mỗi người & Giải thích được luồng \\
\hline
Làm việc nhóm & 1,0 & 1,0 & Thành & Code chung một project \\
\hline
Nâng cao \#3 & 1,5 & 1,5 & Tuấn & Web $\rightarrow$ Buzzer \\
\hline
Nâng cao \#4 & 1,5 & 1,5 & Uy & Lưu theo thời gian \\
\hline
Nâng cao \#5 & 1,5 & 1,5 & Thành & Biểu đồ \\
\hline
Nâng cao \#6 & 1,0 & 1,0 & Tuấn & Telegram \\
\hline
Nâng cao \#9 & 1,5 & 1,5 & Thành & Tài khoản trong database \\
\hline
Nâng cao \#10 & 2,0 & 2,0 & Tuấn (khung) & Tự xây website \\
\hline
Nâng cao \#12 & 2,0 & 1,0 & Uy & WiFiManager \\
\hline
\textbf{Tổng nâng cao} & \textbf{4,0} & \textbf{10,0 $\rightarrow$ 4,0} & & Biên an toàn 6,0 \\
\hline
Báo cáo & 1,0 & 1,0 & Cả nhóm & Nộp đúng hạn \\
\hline
Sản phẩm thật & +1,0 & +1,0 & Uy (lắp ráp) & Cộng vào điểm tổng kết \\
\outerline
\end{tabular}
\end{center}

Hoàn thành toàn bộ, mỗi thành viên đạt $(5 + 4 + 1) \times K = 10 \times K$ điểm đồ án, với $K$ là mức độ hiểu (0--100\%) chấm riêng từng người. Điểm sản phẩm thật được cộng vào điểm tổng kết môn, chỉ xét cho người giải thích được luồng của mình.
\end{minipage}
\par\bigskip

\noindent\begin{minipage}{\linewidth}
\setlength{\parskip}{6pt}
% ======================================================================
\section{Rủi ro khi thực hiện}
% ======================================================================
\begin{tabularx}{\linewidth}{I L{5.0cm} i Y i C{1.6cm} I}
\outerline
\textbf{Rủi ro} & \textbf{Giải pháp} & \textbf{Phụ trách} \\
\hline
MQ-135 cần làm nóng 24--48 giờ & Mua sớm, cấp nguồn để ổn định trước khi hiệu chuẩn & Uy \\
\hline
Hằng số $a$, $b$ và $R_L$ của MQ-135 chưa kiểm chứng & Đối chiếu datasheet gốc; đo $R_L$ bằng đồng hồ vạn năng & Uy \\
\hline
MQ-135 tỏa nhiệt làm sai DHT22 trong hộp & Khoét lỗ thông khí; đặt DHT22 xa MQ-135 & Uy, Thành \\
\hline
LED WS2812 nguồn 5 V nhận tín hiệu 3,3 V & Điện trở 330 $\Omega$; nếu sai màu thì thêm mạch chuyển mức tín hiệu & Tuấn \\
\hline
Ba người cùng sửa firmware trên một board & Mỗi thiết bị một module; quản lý bằng Git; chốt sơ đồ chân từ đầu & Thành \\
\hline
Biểu đồ phụ thuộc dữ liệu lưu & Chốt cấu trúc dữ liệu Firebase ở giai đoạn 1 & Uy, Thành \\
\hline
Backend chạy trên laptop & Laptop bật và có internet suốt buổi demo; sau này chuyển lên VPS & Cả nhóm \\
\hline
Nộp báo cáo trễ thì đồ án bằng 0 & Đặt hạn nội bộ sớm hơn hạn Moodle & Cả nhóm \\
\outerline
\end{tabularx}
\end{minipage}
\par\bigskip

\noindent\begin{minipage}{\linewidth}
\setlength{\parskip}{6pt}
% ======================================================================
\section{Việc cần hỏi giảng viên}
% ======================================================================
\begin{enumerate}
  \item Yêu cầu nâng cao \#10 (tự xây dựng website) có được tính khi website này cũng là nền tảng cho các luồng căn bản không?
  \item Hộp đựng sản phẩm thật cần đáp ứng tiêu chí cụ thể nào để được xét cộng điểm?
  \item Nút BOOT có sẵn trên ESP32 (dùng để xóa cấu hình WiFi) có bị tính là một thiết bị input riêng không?
\end{enumerate}
\end{minipage}
\par\bigskip

\end{document}
